Based on Einstein's general relativity, Russian physicist Alexander Friedmann
derived the Friedmann Equation by which the dynamics of a homogeneous and
isotropic universe can be well described. The Friedmann Equation is usually
written as follows:
\[
\left(\frac{\dot{a}}{a}\right)^2
  = \frac{8\pi G}{3}(\rho_m + \rho_r) + \frac{\Lambda c^2}{3}
    - \frac{kc^2}{a^2}.
\]
We define the Hubble parameter as $H = \dot{a}/a$, where $a$ is the scale
factor and $\dot{a}$ is the rate of change of scale factor with time. Thus,
the Hubble parameter is a function of cosmic time. In the Friedmann Equation,
$\rho_m$ is the density of matter, including dark matter and baryons,
$\rho_r$ is the density of radiation, $\Lambda$ is the cosmological constant,
and $k$ is the curvature of space. Subscript 0 indicates the value of a
physical quantity at present day, e.g.\ $H_0$ is the present value of the
Hubble parameter. Also, to avoid confusion with the reduced Hubble parameter,
we use the reduced Planck constant $\hbar = h/(2\pi)$ instead of the Planck
constant $h$.

\begin{description}
\item[(a)] (5 points) What are the dimensions of the Hubble parameter? One can
  define a characteristic timescale for the expansion of the Universe (i.e.\
  Hubble time $t_H$) using the Hubble parameter. Calculate the present-day
  Hubble time $t_{H0}$.
\item[(b)] (5 points) Let us define the critical density $\rho_c$ as the
  matter density required to explain the expansion of a flat universe without
  any radiation or dark energy. Find an expression for the critical density,
  in terms of $H$ and $G$. Calculate the present critical density
  $\rho_{c0}$.
\item[(c)] (6 points) It is convenient to define all density parameters in a
  dimensionless manner like $\Omega_i = \rho_i/\rho_c$, i.e.\ the ratio of
  density to critical density. The Friedmann Equation can be rewritten using
  these dimensionless density parameters simply as
  \[ \Omega_m + \Omega_r + \Omega_\Lambda + \Omega_k = 1. \]
  Use this information to find expressions for $\Omega_\Lambda$ and
  $\Omega_k$, in terms of $H$, $c$, $\Lambda$, $k$ and $a$.
\item[(d)] (7 points) Another equation which is valid for matter, radiation
  and dark energy is often called the Fluid Equation:
  \[ \dot{\rho} + 3\frac{\dot{a}}{a}\left(\rho + \frac{p}{c^2}\right) = 0, \]
  where $p$ is the pressure of some component, $\rho$ is the density and
  $\dot{\rho}$ is the rate of change of density over time. Radiation contains
  photons and massless neutrinos, and they both travel at the speed of light.
  The pressure exerted by these particles is $1/3$ of their energy density.
  Show that the density of radiation $\rho_r \propto (1+z)^4$, where $z$ is
  cosmological redshift. You may note that if
  $\dot{\rho}/\rho = n\,\dot{a}/a$, then $\rho \propto a^n$.
\item[(e)] (4 points) We know that the value of the cosmological constant
  $\Lambda$ doesn't evolve. Its equation of state has the form
  $p = w\rho_\Lambda c^2$, where $w$ is an integer. Find the value of $w$.
\item[(f)] (13 points) Planck time defines a characteristic timescale before
  which our present physical laws are no longer valid, and where quantum
  gravity is needed. The expression for Planck time can be written in terms of
  $\hbar$, $G$ and $c$, and the non-dimensional coefficient of this expression
  in SI units is of the order of unity. Using dimensional analysis, find an
  expression for the Planck time and estimate its value.
\item[(g)] (7 points) The Planck length defines the length scale associated
  with the Planck time and is given by $l_P = ct_P$. The minimal mass of a
  black hole, also called the Planck mass, is defined as the mass of a black
  hole whose Schwarzschild radius is two times the Planck length. Derive the
  Planck mass $M_P$ and calculate $M_P c^2$ in GeV. This mass is considered to
  be an upper threshold for elementary particles, beyond which they will
  collapse to a black hole.
\item[(h)] (4 points) At the very beginning (soon after the Planck time), all
  the particles were in thermal equilibrium in a primordial soup. As
  temperature decreased, different particles then decoupled from the
  primordial soup one by one and could travel freely in the Universe. Photons
  decoupled at ${\sim}300\,000$ years after the Big Bang. These photons
  emitted at that time are what constitutes the cosmic microwave background
  (CMB), which follows the Stefan--Boltzmann law for blackbody radiation,
  \[ \varepsilon_r = \frac{\pi^2}{15\hbar^3 c^3}(k_B T)^4. \]
  Show that the temperature of the CMB follows $T/(1+z) = \mathrm{constant}$.
\item[(i)] (16 points) With the expansion of the Universe, radiation density
  dropped more quickly than matter density, and at some epoch the matter
  density was equal to the radiation density. Radiation contains both photons
  and neutrinos. Apart from photons, neutrinos additionally contribute to the
  radiation energy density by 68\% (i.e.\
  $\Omega_{r0} = 1.68\,\Omega_{\gamma 0}$, where $\gamma$ indicates photons).
  Estimate the redshift of matter--radiation equality $z_{eq}$ in terms of
  $\Omega_{m0}$ and the reduced Hubble parameter
  $h = H_0/(100\,\mathrm{km\,s^{-1}\,Mpc^{-1}})$. You may use the current
  temperature of the CMB: $T_0 = 2.73$\,K.
\item[(j)] (8 points) The neutrinos decoupled from the primordial soup when
  the temperature of the universe was around 1\,MeV. At this time, the
  radiation density in the universe was much more than all other components.
  Estimate the time ($t = \frac{1}{2H}$) when neutrinos decoupled, and express
  it in seconds since the big bang.
\end{description}
